% 内点、边界点与闭包外的点：A 为开椭圆区域（虚线边界）；y 内点，z 边界点，x 外点。
\begin{tikzpicture}[
  pt/.style={circle,fill=FigureStroke,inner sep=0pt,minimum size=4pt},
  note/.style={font=\FigureNoteFont\fontsize{7.5}{9}\selectfont,text=FigureMuted,align=center}]
  \draw[FigureBlue,fill=FigureBlueFill,line width=.9pt,dashed] (0,0) ellipse (2.2 and 1.3);
  \node[text=FigureBlue] at (-1.25,.45) {$A$};
  % interior point
  \coordinate (y) at (-.3,-.1);
  \draw[FigureGreen,line width=.8pt] (y) circle (.55);
  \node[pt,label={[label distance=-1pt]below:$y$}] at (y){};
  % boundary point at (2.2,0)
  \coordinate (z) at (2.2,0);
  \draw[FigureStroke,line width=.7pt] (z) circle (.42);
  \draw[FigureStroke,line width=.6pt,densely dotted] (z) circle (.22);
  \node[pt,label={[label distance=1pt]above right:$z$}] at (z){};
  % exterior point
  \coordinate (x) at (3.7,1.0);
  \draw[FigureRed,line width=.8pt] (x) circle (.45);
  \node[pt,label={[label distance=-1pt]below:$x$}] at (x){};
  \node[note] at (-.3,-.95) {有球留在$A$内};
  \node[note,anchor=north] at (2.45,-.5) {每个球都跨过边界};
  \node[note,anchor=south] at (3.7,1.5) {有球避开$A$};
\end{tikzpicture}
